\documentclass{ert-l}
\usepackage{amssymb,enumerate}
\overfullrule 5pt
%%\usepackage{showlabels}

%%\textwidth150mm

 %% Operatoren
 %%%%%%%%%%%%%%%%
\newcommand{\Irr}{{\operatorname{Irr}}}

\newcommand{\GL}{{\operatorname{GL}}}
\newcommand{\PSL}{{\operatorname{L}}}
\newcommand{\SL}{{\operatorname{SL}}}
\newcommand{\Sp}{{\operatorname{Sp}}}
\newcommand{\Spin}{{\operatorname{Spin}}}

\newcommand{\tw}[1]{{}^{#1}\!}
\newcommand{\Chevie}{{\sf Chevie}}


 %% theorem style plain --- default
 %%%%%%%%%%%%%%%%%%%%%%%%%%%%%%%%

\newtheorem{thm}{Theorem}%%[section]
\newtheorem{lem}[thm]{Lemma}
\newtheorem{prop}[thm]{Proposition}

\theoremstyle{remark}
\newtheorem{rem}[thm]{Remark}

\raggedbottom


\begin{document}

\title[Corrigendum: Murnaghan--Nakayama rule]
      {Corrigendum: A Murnaghan--Nakayama rule\\ for values of unipotent
       characters in\\ classical groups}

\date{\today}


\author{Frank L\"ubeck}
\address{Lehrstuhl D f\"ur Mathematik, RWTH Aachen, Pontdriesch 14/16,
  52062 Aachen, Germany.}
\email{Frank.Luebeck@math.rwth-aachen.de}

\author{Gunter Malle }
\address{FB Mathematik, TU Kaiserslautern, Postfach 3049,
         67653 Kaisers\-lautern, Germany.}
\email{malle@mathematik.uni-kl.de}

%%\thanks{The second author gratefully acknowledges financial support by ERC
%%  Advanced Grant 291512.}

\keywords{finite simple groups, Cartan invariants}

\subjclass[2010]{Primary 20C20; Secondary 20C33,20C34 }

\begin{abstract}
We settle a missing case in the proof of one of the main applications of
our results in \cite{LM16}.
\end{abstract}

\maketitle

%%%%%%%%%%%%%%%%%%%%%%%%%%%%%%%%%%%%%%%%%%%%%%%%%%%%%%%%%%%%%%%%%%%%%%%%%%%%%
In \cite{LM16} we derived a Murnaghan--Nakayama formula for the values of
unipotent characters of finite classical groups. As one application, we showed
that the first Cartan invariant $c_{11}^{(\ell)}$ of a finite simple classical
group $G$ in non-defining characteristic~$\ell$ is never equal to~2 if Sylow
$\ell$-subgroups of $G$ are
non-cyclic. But contrary to what was claimed in the first line of the proof
of \cite[Thm.~7.1]{LM16}, the same assertion for the case of defining
characteristic does not follow from the work of Koshitani, K\"ulshammer and
Sambale \cite{KKS}. We thank Shigeo Koshitani and J\"urgen M\"uller for
pointing this out to us.

Here, we give a short argument for the missing case. The argument also
covers the
analogous statement in \cite[Prop.~6.3]{LM15} pertaining to simple groups of
exceptional Lie type. The crucial observation is that Sylow $p$-subgroups in
groups of Lie type of characteristic~$p$ are large. Except for a few small
rank cases we do not need any knowledge about the ordinary character table of
the group $G$. 

Note that if $c_{11}^{(p)}=2$ then there exists $\chi\in\Irr(G)$ such that
$1_G+\chi$ is the character of a projective module; in particular $\chi(g)=-1$
for all $p$-singular elements $g\in G$.

\begin{lem}   \label{lem:main}
 Let $G$ be a finite group and $p$ a prime. Assume that $H\le G$ is a
 subgroup such that $|G:H|$ is divisible by $p$, but $|G:H|<|G|_p$. Then
 $c_{11}^{(p)}>2$.
\end{lem}

\begin{proof}
The permutation character $\pi$ of $G$ on the cosets of $H$ has exactly one
trivial constituent $1_G$. On the other hand, as $|G:H|$ is divisible by $p$,
the modular reduction of the permutation representation has (at least) two
trivial modular constituents, one in the socle, one in the head. Thus the
$p$-modular reduction of one of the non-trivial ordinary constituents $\chi$
of $\pi$ must contain the trivial $p$-modular character of $G$. By Brauer
reciprocity,
$\chi$ then occurs in the projective cover $\Psi$ of the trivial character. If
we had $c_{11}^{(p)}=2$ then necessarily $\Psi=1_G+\chi$. Since the degree of
any projective character is divisible by $|G|_p$, this is a contradiction to
$$|G|_p\le\Psi(1)=1_G(1)+\chi(1)\le\pi(1)=|G:H|.$$
\end{proof}

\begin{prop}   \label{main}
 Let $G$ be a finite simple group of Lie type in characteristic~$p$. Then
 $c_{11}^{(p)}>2$ unless $G=\PSL_2(p)$.
\end{prop}

\begin{proof}
If the Lie-rank of $G$ is not very small, then we can find a subgroup of $G$
which fulfills the condition of Lemma~\ref{lem:main}. In Table~\ref{tab} we
give a subgroup $\hat H$ (see e.g.~\cite[\S13]{MT} for the existence) of a
suitable central extension $\hat G$ of $G$.  From the order formulae for 
these groups (see for example~\cite[Tab.~24.1]{MT})
it is easy to check that $\hat H \leq \hat G$ fulfills the condition of
Lemma~\ref{lem:main}. Since the center of $\hat G$ is of order prime to $p$
the same holds for $H \leq G$ where $H$ is the image of $\hat H$ in $G$.

It remains to consider the groups of type $A_1$, $A_2$, $\tw2A_2$,
$\tw3D_4$, $\tw2B_2$, $\tw2G_2$ and $\tw2F_4$.  For all of these the
(generic) character tables are known and available in \Chevie~\cite{Chv}. 
Here we look for all irreducible characters $\chi$ which have value $-1$ on all
$p$-singular (non-semisimple) elements, so that $\chi+1_G$ is a candidate
for the $1$-PIM of $G$.
It is easy to see that apart from two exceptions all characters either
vanish on some $p$-singular element or that a value on a non-trivial
unipotent class is positive or that the values on two non-trivial 
unipotent classes are different.

The two exceptions occur in type $A_1$, so $G = \PSL_2(q)$, which has a
family of characters of degree $q-1$ with value $-1$ on all $p$-singular
elements,
and in type $A_2$, so $G = \PSL_3(q)$, which has a family of characters
of degree $q^3-1$ which also take value $-1$ on all $p$-singular elements.
By \cite[Thm.~5.2]{MW08} the $1$-PIM of $G=\PSL_2(q)$ has degree $(2^f-1)q$,
where $q=p^f$. Since the maximal character degree for $G$ is $q+1$, this
shows that $c_{11}^{(p)}>2$ unless $f=1$, that is, unless $q=p$ is prime, in
which case the Sylow $p$-subgroups of $G$ are cyclic (and so indeed
$c_{11}^{(p)}=2$). 
Similarly, by~\cite[Thm.~5.3]{MW08} the $1$-PIM of $G=\PSL_3(q)$, $q=p^f$,
has degree $(12^f-6^f+1)q^3$, so none of the candidate $1_G+\chi$ mentioned
above can be the $1$-PIM.
\end{proof}

\begin{table}[htbp]
 \caption{Large subgroups}   \label{tab}
\[\begin{array}{|llc||cc|}
\hline
     \hat G&& \hat H&  \hat G& \hat H\cr
\hline
 \SL_4^\pm(q)&& \Sp_4(q)&                   G_2(q)& \tw2A_2(q)\\
 \SL_n^\pm(q)& (n\ge5)& \GL_{n-1}^\pm(q)&   F_4(q)& B_4(q)\\
 \Spin_{2n+1}(q)& (n\ge2)& \Spin_{2n}^-(q)&    E^\pm_6(q)& F_4(q)\\
 \Sp_{2n}(q)& (n\ge3)& \Sp_{2n-2}(q)\Sp_2(q)&      E_7(q)& E_6(q)\\
 \Spin_{2n}^\pm(q)& (n\ge4)& \Spin_{2n-1}(q)&  E_8(q)& E_7(q)\\
\hline
\end{array}\]
\end{table}

Observe that the case $p=2$ was already dealt with in \cite[Thm.~4.5]{KKS}.


%%%%%%%%%%%%%%%%%%%%%%%%%%%%%%%%%%%%%%%%%%%%%%%%%%%%%%%%%%%%%%%%%%%%%%%%%%%%%
\begin{thebibliography}{131}

\bibitem{Chv}
{\sc M.~Geck, G.~Hiss, F.~L\"ubeck, G.~Malle, G.~Pfeiffer},
  {\sf CHEVIE} --- A system for computing and processing generic character
  tables for finite groups of Lie type, Weyl groups and Hecke algebras.
  Appl. Algebra Engrg. Comm. Comput. {\bf7} (1996), 175--210.

\bibitem{KKS}
{\sc S. Koshitani, B. K\"ulshammer, B. Sambale}, On Loewy lengths of blocks.
  Math. Proc. Cambridge Philos. Soc. {\bf156} (2014), 555--570.

\bibitem{LM15}
{\sc C. Lassueur, G. Malle}, Simple endotrivial modules for linear, unitary
  and exceptional groups. Math. Z. {\bf280} (2015), 1047--1074.
  
\bibitem{LM16}
{\sc F. L\"ubeck, G. Malle}, A Murnaghan-Nakayama rule for values of unipotent
  characters in classical groups. Represent. Theory {\bf20} (2016), 139--161.
  
\bibitem{MT}
{\sc G. Malle, D. Testerman}, \emph{Linear Algebraic Groups and Finite Groups
  of Lie Type}. Cambridge University Press, Cambridge, 2011.

\bibitem{MW08}
{\sc G. Malle, Th. Weigel}, Finite groups with minimal 1-PIM. Manuscripta
  Math. {\bf126} (2008), 315--332.
  
\end{thebibliography}

\end{document}
